\section{Coordinate descent and comparisons of compactifications}
\label{sec:equivariant-v169}
We first give an actual gluing statement at a fixed degree. It removes
the dependence on a monic pivot chart without identifying constructions
with different generic Hilbert polynomials.

\begin{theorem}[Equivariant gluing of the fixed-degree graphs]
\label{thm:global-gluing-v169}
Put
\[
 \mathcal P_a=\Pj(H^0(\Pj^1,\OO(a))\otimes k^3),\qquad
 H_a=\Hilb^{P_a}(\Pj^1\times\Pj^2),
\]
and let $\mathcal U_a\subset\mathcal P_a$ be the base-point-free
locus. The closure
\[
 \mathcal G_a=\overline{\{(F,[\operatorname{graph}(F)]):
                                F\in\mathcal U_a\}}
                    \subset\mathcal P_a\times H_a
\]
is an integral projective compactification of the graph of degree-$a$
morphisms with their coefficients retained. It carries a flat universal
embedded curve and a compatible $\PGL_2\times\PGL_3$ action. On every
monic pivot patch it is isomorphic, after the specified undoing of a
target change, to $\A^2\times\Gamma_a$. These patch isomorphisms
satisfy the cocycle identity on all overlaps, including boundary points.
The normalization also carries the action. The morphism
\[
 [\mathcal G_a/(\PGL_2\times\PGL_3)]
       \longrightarrow[\mathcal P_a/(\PGL_2\times\PGL_3)]
\]
is representable and proper, with projective pullback to a scheme
trivializing the group torsor.
\end{theorem}
\begin{proof}
The Hilbert scheme is projective and separated, so the closure is
projective and integral and is an isomorphism over $\mathcal U_a$.
Pullback of the universal Hilbert family is flat. A pair of source and
target automorphisms takes a graph to the graph of the transformed
map, preserves $\OO(1,1)$, and hence acts on its Hilbert scheme.
It preserves the closure and its universal curve.

Here are the patch maps explicitly. Choose a source point where some
coordinate form is nonzero, send it to $[0:1]$, and permute that target
coordinate into position zero. In the projective coefficient patch
normalize its $t^a$ coefficient to one and write
\[
 F_0=f^{\mathrm h},\qquad
 F_1=\alpha_1 f^{\mathrm h}+g^{\mathrm h},\qquad
 F_2=\alpha_2 f^{\mathrm h}+r^{\mathrm h},
 \qquad\deg g,\deg r<a.
\]
This identifies the coefficient patch with $\A^2\times B_a$.
The target matrix and its inverse are
\begin{equation}\label{eq:undo-global-v169}
 U(\alpha)=
 \begin{pmatrix}1&0&0\\-\alpha_1&1&0\\-\alpha_2&0&1\end{pmatrix},
 \qquad
 U(\alpha)^{-1}=
 \begin{pmatrix}1&0&0\\\alpha_1&1&0\\\alpha_2&0&1\end{pmatrix}.
\end{equation}
The first matrix identifies the graph with the polynomial-contact
graph; the second returns it to the original fixed target. Since the
product with $\A^2$ is flat, taking the closure of the generic graph
commutes with this product. Thus the closure on the patch is exactly
the asserted product, with its scheme structure.

The patches cover $\mathcal P_a$: a nonzero form cannot vanish at
$a+1$ distinct chosen source points, so finitely many such choices
and the three target pivots suffice. On an overlap the change is the
composition of the actual source matrices, the target permutations,
and \eqref{eq:undo-global-v169}. Its inverse reverses that composition.
The product of the matrices on a triple overlap is the identity as a
projective transformation. The induced Hilbert maps therefore satisfy
the same identity. No comparison based only on reduced closed points
is used.

For normalization, the product of the smooth connected group with
the normal variety $\mathcal G_a^\nu$ is normal. The action composed
with normalization is dominant and lifts uniquely to the normalization;
uniqueness proves the action identities. Finally an equivariant
morphism induces a representable map of the quotient stacks. After a
torsor is trivialized its pullback is the projective morphism
$\mathcal G_a\to\mathcal P_a$; properness descends.
\end{proof}
This theorem glues \emph{all} pivot charts of a fixed-degree rational-map
problem. Its construction is classical graph compactification; the
explicit patch maps state exactly how the contact graphs sit in it.
It is not a gluing of every Smith and Kronecker stratum of every
quadratic-pencil Hilbert problem. In particular, an adjugate tuple that
vanishes identically supplies no point of $\mathcal P_a$.

\begin{corollary}[Moving contacts and mixed target coefficients]
\label{cor:moving-v169}
Over $\A^1_\alpha\times S_a$ the family
\[
 ((x-\alpha)^a,b,c(x-\alpha)^{a-1})
\]
has strict Hilbert graph $\A^1_\alpha\times T_a$, with the universal
curve obtained by substituting $x-\alpha$ for $x$. Over the further
open parameter space $1-uv\ne0$, the same statement holds for
\begin{equation}\label{eq:mixed-family-v169}
 ((x-\alpha)^a,\ b+uc(x-\alpha)^{a-1},\ vb+c(x-\alpha)^{a-1}),
\end{equation}
after undoing the target matrix
$\left(\begin{smallmatrix}1&0&0\\0&1&u\\0&v&1\end{smallmatrix}\right)$.
For $a=2$ both latter coordinate polynomials can have nonzero constant
and linear terms. This family is an equivariant enlargement of the
marked slice, not the unrestricted six-coefficient contact family.
\end{corollary}
\begin{proof}
Translation is a source automorphism over $\A^1_\alpha$. The indicated
target matrix is invertible precisely on the stated open, with inverse
lower block $(1-uv)^{-1}\left(\begin{smallmatrix}1&-u\\-v&1\end{smallmatrix}\right)$.
They identify the entire ambient families, not just their generic
fibres. Applying these isomorphisms to the flat universal curve and
its graph closure proves the result.
\end{proof}

\begin{proposition}[Conservation of the Hilbert polynomial]
\label{prop:conservation-v169}
A projective flat finitely presented family over a connected base
cannot have, as its fibres, just the graphs of morphisms of two
different degrees. A degree drop from $a$ to $a-d$ in a limit of
$\mathcal G_a$ must be accompanied by additional scheme structure
whose Hilbert polynomial contributes the missing degree and Euler
characteristic. The extra structure may include punctual nilpotents,
as in Theorem~\ref{thm:genus-correction-v169}.
\end{proposition}
\begin{proof}
The Hilbert polynomial for $\OO(1,1)$ is locally constant in a
projective flat family, whereas the two graphs have polynomials
$(a+1)l+1$ and $(a-d+1)l+1$. These differ for $d>0$. In a fixed-degree
Hilbert compactification the polynomial is retained by the whole
limiting scheme, not by deleting its tails or embedded associated
points.
\end{proof}
Thus comparison across different generic polynomials cannot consist
of an equality of the unaugmented flat graph functors. The horizontal
comparison and common-refinement results in
Section~\ref{sec:comparison-v167} continue to apply on a fixed dense
open with a fixed generic embedded quotient. They are not changed into
a false full-pullback statement by the equivariant gluing above.

\subsection{Tropical data, residues, and coordinate changes}
\begin{proposition}[Comparison of coefficient presentations]
\label{prop:fan-comparison-v169}
For a finite list of polynomial projective coordinates, the finite
weight-layer and residue construction of
Theorem~\ref{thm:finite-fan-v167} extends to exact monomial profiles
with nonnegative weights, including specified unit coordinates, and
to polynomial coordinates translated about any fixed centre. It gives
finitely many algebraic families of closed projective points, not
finitely many embedded curves or abstract isomorphism types.

Under an invertible monomial change of coefficient torus coordinates,
the weights transform by the exponent matrix and the residues by the
same monomial substitution; the finite partitions have a common
pullback refinement. Under a general polynomial coordinate change,
the corresponding statement holds after composing the projective
coordinate polynomials with that change. An exact monomial arc need
not remain monomial in the new coordinates, and the two displayed
fans need not be isomorphic.
\end{proposition}
\begin{proof}
Expand each of the finitely many polynomials into its finite monomial
support, after any specified translation or polynomial substitution.
The equalities of weights partition the nonnegative orthant into a
finite rational polyhedral arrangement. Within a cone, group terms
of equal weight and partition the residue torus by the first nonzero
coefficient sum in each polynomial. Zero polynomials are kept as zero;
identically zero input coordinates are a separate designated pattern.
Every choice is finite, and the residual vector is polynomial on each
locally closed piece. Weight zero introduces constant residue terms
but does not change the argument. For a monomial substitution the
weight and residue formulas follow directly from substitution.
For a polynomial substitution repeat the finite expansion on the
original monomial arc; this compares the resulting points even when
there is no transformation of fans by a linear weight map. For
localized-polynomial changes work on an open where denominators are
units and clear them by a common denominator, which does not change
the projective map.
\end{proof}

Applied to maximal minors, valuations are realizable valuated-matroid
data: the minors satisfy all Pl\"ucker equations, and the minimum
valuation in each nonzero relation is attained at least twice. On the coordinate stratum of the nonzero minors these are points
of the tropicalization of that Grassmannian stratum, as in
Speyer--Sturmfels \cite[Theorem~2.1 and Section~3]{SpeyerSturmfels169}.
It is not enough to check only a selected set of tropical quadratic
relations and then claim realizability. Valuations forget residues;
our first wall already has the same valuations for every $\kappa\ne0$
but different embedded ideals. The primitive Hilbert vector retains
these residues and determines the degree-$m$ kernel. It does not
select among normalization branches over a graph point.

On $S_a$ the individual nonzero minors are monomials. The normal fan
of their Newton polygon is exactly the chain fan, after removal of
content; all its walls are effective. This identifies the computed fan
with the restriction to nonnegative weights of the state fan of this marked torus orbit, in the classical
initial-ideal setting of \cite{BayerMorrison167}. No such minimality
claim is made for the full coefficient-support arrangement. In
comprehensive Gr\"obner terminology, which organizes specializations
by leading-coefficient conditions \cite{ManubensMontes169},
Proposition~\ref{prop:curve-charts-v169} gives a particularly simple
system: a single monic basis on every affine parameter chart, with no
further coefficient conditions. Its explicit standard monomials,
rather than a Hilbert-scheme elimination, recover the curve ideals.

\begin{proposition}[Equality up to a target automorphism]
\label{prop:target-equivalence-v169}
Let two embedded curves with polynomial $P_a$ have degree-$m$
ideal kernels $K,K'\subset H^0(X,\OO(m,m))$, $m\geq m_a$.
They differ by a target automorphism exactly when there is
$A\in\PGL_3$ for which
\[
 (1\otimes\operatorname{Sym}^m A^*)K'=K.
\]
This is an invariant, finite algebraic criterion: equality of the
two fixed-rank subspaces is given by rank conditions, with $\det A\ne0$.
It does not classify abstract curve isomorphisms. On a fixed wall
of the marked family, all nonzero residue values are related by a
diagonal target automorphism, although their fixed-target Hilbert
points differ.
\end{proposition}
\begin{proof}
An automorphism acts on sections by pullback. The Hilbert immersion
in degree $m$ makes equality of the kernels equivalent to equality
of ideal sheaves. This proves both implications and the algebraic
criterion. Scaling $F,G,R$ independently scales $\kappa$, or scales
$\lambda$ by the ratio between $R^j$ and $F^{j-1}G$; every nonzero
ratio is available over $k$. No assertion about different walls is
needed for this last observation.
\end{proof}

\subsection{A concrete comparison with Quot and stable-map limits}
\begin{proposition}[A Quot contraction of the entire exceptional chain]
\label{prop:quot-contraction-v169}
The marked family defines a flat quotient on the fixed source
\begin{equation}\label{eq:quot-sequence-v169}
 0\longrightarrow\OO_{\Pj^1}(-a)
   \xrightarrow{(t^a,bs^a,cst^{a-1})}\OO_{\Pj^1}^{\oplus3}
   \longrightarrow\mathcal Q\longrightarrow0
\end{equation}
over all of $S_a$. The resulting morphism from $T_a$ to the fixed-source
Quot scheme is constant on its entire exceptional chain. The common
central quotient is
\[
 \mathcal Q_0\simeq\OO_{a[1:0]}\oplus\OO_{\Pj^1}^{\oplus2}.
\]
After adding three fixed source markings away from $[1:0]$, this also
gives stable quotients in the sense of Marian--Oprea--Pandharipande.
In contrast, the Hilbert map distinguishes every point of the chain.
\end{proposition}
\begin{proof}
The first coordinate $t^a$ makes the sheaf map fibrewise injective,
even when $b=c=0$. Both source and target bundles are flat over $S_a$;
the fibrewise injectivity criterion for this finite-presentation
sequence gives flatness of its cokernel over $S_a$. Thus it defines
a morphism to the Quot scheme already before the blow-up. At the
origin the sequence is multiplication by $t^a$ in its first summand,
which gives the stated quotient. Its pullback along $T_a\to S_a$
therefore contracts the chain. The Hilbert map on that chain is a
closed immersion because $T_a$ is a closed graph in $S_a\times H_a$.
For three markings the source is stable, the torsion is away from
the markings, and $\omega_{\Pj^1}(p_1+p_2+p_3)\otimes\OO(a\epsilon)$
has positive degree for every $\epsilon>0$. These verify the actual
stability conditions of \cite[Section~2]{MOP169}; no unstable
unmarked rational source is being called a stable quotient.
\end{proof}

\begin{corollary}[Non-Cohen--Macaulay limits of smooth rational graphs]
\label{cor:external-curves-v169}
For each $a\geq3$ the Hilbert closure of smooth rational graphs in
$\Pj^1\times\Pj^2$ with polynomial $P_a$ contains a curve whose
reduction is a line attached to a rational cuspidal plane curve of
degree $a$, and whose punctual nilradical has length $g_a$ and
nilpotence order $a-1$. This curve is invisible as a distinct punctual
structure to the associated-cycle map. It is not the scheme-theoretic
image of a map from a reduced nodal curve.
\end{corollary}
\begin{proof}
Take $b=\tau^a,c=\tau$. The generic graph is smooth because its
projection to the source is an isomorphism. Its limit is
$\mathcal W_a$ with $\lambda=1$, so
Theorem~\ref{thm:genus-correction-v169} proves the first assertions.
An associated one-cycle records generic lengths at one-dimensional
components and discards finite-length submodules. Finally the
scheme-theoretic image of a morphism from a reduced scheme is reduced:
if a local target function has a power mapping to zero, its pullback
is nilpotent and hence zero. The asserted nonreduced curve cannot
be such an image. This does not rule out comparison spaces between
stable-map and Hilbert compactifications; it rules out identifying
this Hilbert limit with the unmodified scheme image of a stable map.
\end{proof}
This application is independent of the finite failure algebra and of
Paper I. General comparisons between stable maps and stable quotients
already involve spaces dominating both constructions
\cite{Manolache169}; the explicit contraction above is a different,
fixed-source Hilbert-to-Quot calculation, not a replacement for those
virtual-class comparison theorems.

\subsection{Return to the reconstructed pencil}
A specified effective failure family first reconstructs its pencil by
Paper I. Wherever its fixed-target contact comparison lands in $B_a$,
the graph in Theorem~\ref{thm:global-gluing-v169} supplies compatible
models under source and target frame changes. If its coefficient map
factors through \eqref{eq:marked-family-v169} or
\eqref{eq:mixed-family-v169}, the chain, all local rings, and the
punctual modules above pull back as strict horizontal families.
The inherited complete-quadric comparison determines the target in
that application; forgetting that target and its coefficient map
would lose the meaning of the embedded limit. This statement does
not identify every complete-quadric Hilbert component with
$\mathcal G_a$, does not cover an identically zero adjugate tuple,
and does not turn Kronecker data alone into a higher-corank fibre
classification. No boundary operation on an isolated unframed
multiplication table is inferred.
